\begin{topic}[id=p5g_wifi7_robot_iiot,
             difficulty={Anspruchsvoll},
             type={Vergleichs- und Netzarchitekturanalyse},
             prerequisites={Grundlagen zu Rechnernetzen, industrieller Kommunikation, Funknetzen und IT-Sicherheit},
             practice={gut möglich},
             owner={Dozententeam},
             updated={2026-04-13},
             tags={ Private 5G, Wi-Fi 7, mobile Roboter, IIoT, Campusnetze, Roaming, QoS, Netzplanung, industrielle Funknetze, Energieeffizienz }]{Private 5G vs. Wi-Fi 7 für mobile Roboter und IIoT}

\TopicDescription{Private 5G und Wi-Fi 7 konkurrieren in Fabriken, Logistikflächen und Campusnetzen als drahtlose Infrastruktur für mobile Roboter und IIoT-Geräte, unterscheiden sich aber grundlegend in Architektur, Betriebsmodell und Planungslogik. Für ein seminarfähiges Thema steht daher nicht die Frage nach maximaler Datenrate im Vordergrund, sondern die systematische Bewertung von Latenz, Roaming-Verhalten, QoS-Durchsetzung, Energiebedarf und Funkplanung unter realen Mobilitätsanforderungen. Analysiert wird, wie Private 5G mit Non-Public-Network-Architekturen, zentralem Mobilitätsmanagement, priorisierten QoS-Flows, SIM-basierter Authentisierung und kontrollierter Campusabdeckung arbeitet und wie sich dies von Wi-Fi-7-Netzen mit 802.11be, Multi-Link-Betrieb, WPA3-basierter Absicherung, hohem Durchsatz und typischer Enterprise-Planung unterscheidet. Relevant sind dabei typische Einsatzfälle wie AMRs, fahrerlose Transportsysteme, Teleoperation, mobile Sensorik und gemischt drahtlose Produktionszellen. Einbezogen wird außerdem, wie Netzplanung, Standortdesign und Störungsreserven die tatsächliche Nutzbarkeit im Feld beeinflussen. Ebenso wichtig sind die Grenzen beider Ansätze: 5G erfordert Spektrum, Kernnetz- und SIM-Management sowie sorgfältige Integration in OT- und IT-Landschaften, während Wi-Fi 7 bei dichter Belegung, unterbrechungsarmem Handover und streng garantierbarer Dienstqualität an systemische Grenzen stoßen kann. Das Thema eignet sich deshalb für eine vergleichende Architektur- und Engineeringanalyse zwischen Flexibilität, Planbarkeit, Betriebsaufwand, Sicherheitsintegration und Interoperabilität in industriellen Campusnetzen.}

\TopicLeitfragen{%
\item Welche Mobilitäts- und QoS-Anforderungen sprechen in Campusnetzen eher für Private 5G, welche eher für Wi-Fi 7?
\item Wie unterscheiden sich Roaming, Netzplanung und Funkzellenkonzepte für AMRs und mobile IIoT-Geräte?
\item Welche Grenzen bestehen bei deterministischer Latenz, Energiebedarf und Interferenzrobustheit beider Ansätze?
\item Wie wirken sich Betriebsmodell, Sicherheitsarchitektur und Engineering-Aufwand auf die langfristige Skalierbarkeit aus?
}

\TopicScope{%
\item Keine vollständige Marktübersicht zu kommerziellen Access Points, 5G-Kernnetzen oder Campusnetz-Anbietern.
\item Kein detaillierter Frequenzrechtsvergleich nationaler Spektrum- und Lizenzmodelle.
\item Keine experimentelle Performanzmessung mit realer Funkhardware im Seminarumfang.
\item Kein Fokus auf 6G, LoRaWAN oder Bluetooth-basierte Nebenlösungen.
}

\TopicPractice{Für ein Beispiel mit AMRs in einer Produktionshalle und mobiler IIoT-Sensorik in einem Campusnetz kann eine Vergleichsmatrix mit Funkzellen, Übergabepfaden, QoS-Klassen, Energieprofilen und Betriebsaufwand erstellt werden. Dazu reichen Architekturdiagramme und ein strukturierter Engineering-Vergleich ohne Hardware-Aufbau.}

\TopicKeywords{Private 5G, Wi-Fi 7, mobile Roboter, IIoT, Campusnetze, Roaming, QoS, Netzplanung, industrielle Funknetze, Energieeffizienz}

\nocite{net_p5gwifi7_3gpp_rel18,net_p5gwifi7_3gpp_npn,net_p5gwifi7_acia_npn2024,net_p5gwifi7_acia_devices2022,net_p5gwifi7_ieee80211be2024,net_p5gwifi7_wfa_cert2023}
\end{topic}
